\section{Motivación y contexto del problema}
\label{sec:motivacion}

\subsection{El sistema estudiado}
Una empresa de logística, un servicio de emergencias o un algoritmo de \textit{clustering}
enfrentan la misma decisión esencial: \emph{¿dónde instalar un número fijo de centros para
atender, lo más cerca posible, a un conjunto de puntos de demanda?} Pensemos en una ciudad con
$N$ barrios y $M$ ubicaciones candidatas para bodegas; se dispone de presupuesto para abrir
exactamente $p$ bodegas; cada barrio será servido por la bodega abierta más cercana. La pregunta
es qué $p$ ubicaciones minimizan la suma de distancias barrio--bodega. Este es el
\textbf{problema de las $p$-medianas} (pMP).

\begin{intuicion}
``Pongo $p$ antenas; cada casa se conecta a la antena más cercana; quiero minimizar el largo
total de cable.'' Lo notable es que, \emph{una vez fijadas las antenas}, cada casa decide
trivialmente (elige la más cercana). Lo difícil es decidir \emph{cuáles} $p$ antenas.
\end{intuicion}

\todo{FALTA: figura/esquema del sistema (clientes y sitios candidatos, asignación al sitio
abierto más cercano), idealmente sobre la instancia \texttt{toy1}.}

\paragraph{Naturaleza de este trabajo.} Este informe es una \textbf{replicación y validación}
del método de Duran-Mateluna, Ales \& Elloumi (2023). El hilo conductor es la teoría del curso
—estructura, relajación, dualidad, cotas, descomposición, algoritmo— aplicada a un problema
concreto; la implementación en C es \emph{evidencia} de que la teoría funciona, no la
contribución en sí. Las familias de instancias (OR-Library, TSP-Library) aparecen porque son los
\emph{benchmarks estándar} que usa el paper, no como eje del trabajo.

\subsection{Relevancia}
El pMP es un problema fundamental de localización discreta, sin costos fijos de apertura (lo que
lo distingue del problema de localización de instalaciones, UFL) y con el número de aperturas
\emph{fijado} en $p$. Aparece en:
\begin{itemize}[noitemsep]
  \item \textbf{Logística y cadena de suministro:} ubicación de bodegas, plantas, centros de distribución.
  \item \textbf{Respuesta a emergencias y ayuda humanitaria:} ubicación de refugios o puntos de acopio.
  \item \textbf{Aprendizaje no supervisado:} el \textit{k-medoids}, donde clientes y sitios son el
        mismo conjunto de puntos y se buscan $p$ ``centros'' representativos.
\end{itemize}

\subsection{Por qué es computacionalmente difícil}
El pMP es \textbf{NP-duro} (Kariv \& Hakimi, 1979). El número de formas de elegir $p$ sitios entre
$M$ es $\binom{M}{p}$, que crece de manera combinatoria. Las formulaciones MILP exactas crecen
rápido en variables y restricciones; para instancias con decenas o cientos de miles de puntos los
solvers genéricos colapsan en memoria o tiempo. El desafío abierto que ataca el paper de
referencia es resolver \emph{exactamente} a gran escala.

\subsection{Qué aporta el método replicado}
Duran-Mateluna, Ales \& Elloumi (2023) aplican \textbf{descomposición de Benders} a la mejor
formulación exacta conocida. El hallazgo clave es que el subproblema de Benders del pMP es tan
simple que su \emph{dual tiene solución cerrada}: los cortes se separan en tiempo lineal $O(NM)$
\textbf{sin resolver ningún programa lineal}. Esto, dentro de un esquema de dos fases (relajación
+ branch-and-cut con cortes perezosos), supera al método previo estado-del-arte (Zebra) por un
orden de magnitud y resuelve por primera vez a optimalidad instancias enormes.

Este informe integra las tres entregas teóricas del curso (formulación, relajación lagrangiana,
descomposición de Benders) y reporta una \textbf{implementación propia en C} con resultados
reales. El detalle teórico extendido está en \texttt{docs/PMEDIAN\_BENDERS\_PROJECT\_BRIEF.md}.
